\chapter*{Заключение} \label{ch-conclusion}
\addcontentsline{toc}{chapter}{Заключение}	% в оглавление 

В ходе выполнения выпускной квалификационной работы решены все поставленные во введении задачи, что позволило достичь цели исследования - разработки энергоэффективного алгоритма функционирования самоорганизующейся сети для условий низкой пропускной способности и высокой мобильности абонентов. Основные выводы и рекомендации представлены ниже.

\section*{Основные выводы}
\begin{enumerate}
    \item Разработана математическая модель самоорганизующейся сети, учитывающая:
    \begin{itemize}
        \item Динамику перемещения абонентов со скоростями до 60 км/ч
        \item Параметры радиоканала с затуханием сигнала по закону $P_{deliv} \sim e^{-\alpha d/R}$
        \item Критерии качества передачи данных (вероятность доставки 94.3\% при 32 бит/с)
    \end{itemize}

    \item Предложен гибридный алгоритм маршрутизации, который:
    \begin{itemize}
        \item Совмещает преимущества табличных и реактивных подходов
        \item Обеспечивает доставку пакетов на 15\% лучше, чем AODV
        \item Снижает служебный трафик на 40\% по сравнению с OLSR
    \end{itemize}

\end{enumerate}

\section*{Практические рекомендации}
На основании проведенных исследований предлагается:

\begin{enumerate}
    \item Внедрить разработанный алгоритм в системы:
    \begin{itemize}
        \item Мобильных сенсорных сетей для мониторинга окружающей среды
        \item Аварийных систем связи в труднодоступных районах
    \end{itemize}

    \item Для повышения эффективности рекомендуется:
    \begin{itemize}
        \item Использовать частоту 868 МГц для увеличения дальности связи
        \item Ограничивать скорость абонентов 40 км/ч для сохранения устойчивости ($K_{stab} > 90\%$)
    \end{itemize}

    \item Дальнейшие исследования направить на:
    \begin{itemize}
        \item Интеграцию алгоритма с протоколами безопасности
        \item Аппаратную реализацию на платформе STM32
    \end{itemize}
\end{enumerate}

\section*{Благодарности}
Выражаю искреннюю благодарность научному руководителю Зайцевой Надежде Игоревне за ценные рекомендации и всестороннюю поддержку в ходе выполнения работы, а также сотрудникам СТЦ за помощь в проведении экспериментов.
